\documentclass[italian,12pt]{article} %tipo di documento

\def\Title{Verbale interno}
\def\Subtitle{Riunione interna settimanale}
\def\Author{7Last}
\def\Date{2024-05-22}
\def\Version{v1.0}

%pacchetti extra da scaricare dblfloatfix, fancyhdr
\usepackage[left=2cm, right=2cm, bottom=3cm, top=3cm]{geometry}
\usepackage{fancyhdr}%creazione header-footer
\usepackage{graphicx} %serve per inserire immagini
\graphicspath{ {../../../logo/} }
%\usepackage{dblfloatfix} %serve per posizionare gli elementi dove si vuole
\usepackage[hidelinks]{hyperref} %serve per i link
\usepackage{tikz}
\usepackage{tgadventor} % font
\usepackage[useregional=numeric,showseconds=true,showzone=false]{datetime2}
\usepackage{caption}

\usepackage{hyperref}
\usepackage{tocloft}
\usepackage{titlesec}
\usepackage{color}
\usepackage{ulem}
\usepackage{pgfplots}
\usepackage{pgf-pie}
\usepackage[italian]{babel}
\usepackage{comment}
\usepackage{tabularx}
\usepackage{longtable}
\usepackage{float}
\usepackage{amsmath}

\linespread{1.2}
\captionsetup[table]{labelformat=empty}
\geometry{headsep=1.5cm}

\renewcommand{\contentsname}{Indice}%imposto il nome dell'indice
\renewcommand\familydefault{\sfdefault}

%-------------------INIZIO DOCUMENTO--------------
\begin{document}

\newgeometry{left=2cm,right=2cm,bottom=2.1cm,top=2.1cm}
\begin{titlepage}
    \vspace*{.5cm}

    \vspace{2cm}
    {
        \centering
        {\bfseries\huge \Title\par}
        \bigbreak
        {\bfseries\large \Author\par}
        \bigbreak
        {\Version\par}
        \vfill

        \begin{tikzpicture}[remember picture,overlay]

            \fill[blue!20!black] (current page.south west) -- ++(10cm,0) -- ++(-10cm,15cm) -- cycle;
            \fill[orange] (current page.south east) -- ++(-18cm,0) -- ++(21.6cm,22cm) -- cycle;

            \clip (0,-3cm) circle (2.5cm) node (current page.center) {\includegraphics[width=5cm]{logo.jpg}};
        \end{tikzpicture}

    }

    \vfill

\end{titlepage}

\restoregeometry






















\newpage

%---------------header------------------%
\pagestyle{fancy}
\fancyhead{} % pulizia degli header
\lhead{%
\begin{tikzpicture}
    \clip (0,0) circle (0.5cm);
    \node at (0,0) {\includegraphics[width=1cm]{logo}};
\end{tikzpicture}%
}
\chead{\vspace{\fill}\Title\vspace{\fill}}
\rhead{\vspace{\fill}\Version\vspace{\fill}}
%quad è una spaziatura
%---------------------------------------%

%-----------tabella revisioni-----------%
\begin{table}[!h]
	\caption{Versioni}
	\footnotesize
	\begin{center}
		\begin{tabular}{ l l l l l }
			\hline                                                                          \\[-2ex]
			Ver. & Data       & Redattore         & Verificatore       & Descrizione        \\
			\\[-2ex] \hline \\[-1.5ex]
			1.0  & 2024-04-09 & Antonio Benetazzo & Valerio Occhinegro & Stesura definitiva \\
			0.1  & 2024-03-28 & Antonio Benetazzo & Valerio Occhinegro & Prima stesura      \\
			\\[-1.5ex] \hline
		\end{tabular}
	\end{center}
\end{table}
%---------------------------------------%

\newpage

\tableofcontents

\newpage

\section{Dettagli della riunione}


\textbf{Sede della riunione}: Piattaforma Discord\\
\textbf{Orario di inizio}: 21:00\\
\textbf{Orario di fine}: 23:00\\


\begin{flushleft}
	\begin{table}[!h]
		\begin{tabular}{ |l|l|l| }
			\hline
			\textbf{Partecipante} & \textbf{Ruolo} & \textbf{Presenza} \\
			\hline
			Antonio Benetazzo     & Redattore      & Presente          \\
			Davide Malgarise      &                & Assente           \\
			Elena Ferro           &                & Presente          \\
			Leonardo Baldo        &                & Presente          \\
			Matteo Tiozzo         & Amministratore & Presente          \\
			Raul Seganfreddo      &                & Presente          \\
			Valerio Occhinegro    & Verificatore   & Presente          \\
			\hline
		\end{tabular}
	\end{table}
\end{flushleft}

\section{Ordine del giorno}
\begin{itemize}
	\itemsep0em
	\item Approvazione della candidatura
	\item Prossimi step e documentazione da produrre
	\item Aggiornamento \textit{Way of Working}
	\item Modifiche al template per la documentazione
	\item Varie ed eventuali
\end{itemize}

\newpage

\section{Verbale}

\subsection{Approvazione della candidatura}
Prendiamo atto dell'approvazione della candidatura da noi presentata per il progetto \textit{SyncCity} da parte del professor Vardanega. \\
Nella valutazione viene evidenziato come la pianificazione sia ancora molto \textit{ingenua}. Discutiamo assieme su tali osservazioni e ci riserviamo di rivalutare tale pianificazione in occasione della pubblicazione della prossima documentazione, prevista per il prossimo periodo. Ci viene inoltre suggerito di non tardare nell'impostazione del documento \textit{Piano di Progetto} di cui discuteremo in uno dei punti successivi. \\
Siamo comunque soddisfatti del risultato raggiunto e ci impegniamo a migliorare la pianificazione per le prossime scadenze. \\
Provvediamo anche a comunicare via e-mail a \textit{Sync Lab} l'approvazione della candidatura e alle altre due aziende proponenti gli altri capitolati la nostra rinuncia.

\subsection{Prossimi step e documentazione da produrre}
Discutiamo insieme i prossimi passi: le prossime revisioni saranno la RTB (Requirements and Technology Baseline), la PB (Product Baseline) e l'eventuale CA (Customer Acceptance).\\
Esaminiamo la documentazione necessaria per la prossima revisione, l'RTB (che utilizzeremo per il resto del progetto): \textit{Analisi dei Requisiti}, \textit{Piano di Progetto}, \textit{Piano di Qualifica}, \textit{Norme di Progetto}, \textit{Glossario},  oltre alla \textit{Lettera di Presentazione}, al \textit{POC} (Proof of Concept) e a tutti i verbali interni ed esterni. Ci impegniamo a studiare questi documenti nel prossimo periodo per comprenderne la struttura e i contenuti attesi e ad iniziare a lavorarci il prima possibile. \\
In particolare, per quanto riguarda il Glossario, riteniamo utile trovare un modo per evidenziare automaticamente ogni termine in tutti i documenti in cui compare, in modo che sia chiaro al lettore che si tratta di un termine incluso in esso, senza doverlo gestire manualmente. \\
Riteniamo inoltre opportuno fissare una chiamata con l'azienda per discutere come affrontare le prossime fasi del progetto.

\subsection{Aggiornamento \textit{Way of Working}}
Risulta necessario e sempre più urgente anche procedere con l'aggiornamento del \textit{Way of Working} per la pianificazione e la gestione del lavoro nelle fasi successive. \\
In particolare, nei prossimi giorni valuteremo insieme l'utilizzo di \textit{ClickUp} come strumento di gestione del lavoro e delle scadenze, come suggeritoci dall'azienda proponente. Valuteremo anche eventuali altri strumenti a supporto della gestione del progetto. Provvederemo inoltre, come già detto, alla stesura del documento \textit{Norme di Progetto}, dove potremo definire meglio le modalità di lavoro e le regole da seguire.

\subsection{Modifiche al template per la documentazione}
Decidiamo di apportare alcune modifiche al template per la documentazione. \\
Oltre ad alcuni accorgimenti grafici, valutiamo sia il caso di dividerlo in tre differenti template: uno per la documentazione in genere, uno specifico per i verbali interni (che riporteranno anche la data, il registro delle presenze, l'ordine del giorno e una struttura definita per il verbale) e uno specifico per i verbali esterni (che conterrà in aggiunta la struttura necessaria per la firma da parte dell'azienda proponente).

\subsection{Conclusioni}
Ci diamo appuntamento alla prossima settimana, riteniamo sia importante trovarci possibilmente almeno una volta a settimana per valutare assieme l'andamento del progetto e per discutere eventuali problematiche o dubbi.

\end{document}



